\documentclass[11pt,a4paper]{article}

\usepackage[utf8]{inputenc}
\usepackage[T1]{fontenc}
\usepackage[brazil]{babel}
\usepackage{lmodern}
\usepackage{geometry}
\usepackage{microtype}
\usepackage{xcolor}
\usepackage{graphicx}
\usepackage{tabularx}
\usepackage{booktabs}
\usepackage{enumitem}
\usepackage{tikz}
\usepackage{fancyhdr}
\usepackage{hyperref}

\usetikzlibrary{arrows.meta,positioning}
\geometry{margin=2.3cm, headheight=14pt}
\setlength{\parindent}{0pt}
\setlength{\parskip}{0.6em}
\setlist[itemize]{leftmargin=1.5em,itemsep=0.2em}
\setlist[enumerate]{leftmargin=1.5em,itemsep=0.2em}

% ---------------------------------------------------------------
% Dados do projeto: altere aqui e o resto do documento acompanha
% ---------------------------------------------------------------
\newcommand{\jogo}{The Last Contract}
\newcommand{\subtitulo}{Mini Game Design Document}
\newcommand{\equipe}{André Bartzen Araujo \and Jackson Matiazzi \and José Maciel Camargo}
\newcommand{\disciplina}{Disciplina --- Avaliação Grau 2}
\newcommand{\gddversao}{0.36}

\definecolor{Destaque}{HTML}{7A2E2E}
\definecolor{Pendente}{HTML}{9A9A9A}

% Texto ainda não preenchido aparece em cinza e itálico.
% Para achar o que falta: buscar por \preencher no .tex
\newcommand{\preencher}[1]{\textcolor{Pendente}{\textit{[#1]}}}

\pagestyle{fancy}
\fancyhf{}
\fancyhead[L]{\textcolor{Destaque}{\jogo}}
\fancyhead[R]{\subtitulo{} --- v\gddversao}
\fancyfoot[C]{\thepage}
\renewcommand{\headrulewidth}{0.4pt}

\hypersetup{
    colorlinks=true,
    linkcolor=black,
    urlcolor=Destaque,
    pdftitle={\jogo{} - \subtitulo}
}

\title{\textbf{\jogo} \\ \large \subtitulo{} --- versão \gddversao}
\author{\equipe}
\date{\disciplina \\ \today}

\begin{document}

\maketitle
\thispagestyle{empty}
\tableofcontents
\newpage

% ===============================================================
\section{Ficha técnica}

\begin{tabularx}{\textwidth}{@{}lX@{}}
\toprule
\textbf{Título}          & \jogo \\
\textbf{Gênero}          & Ação 3D \\
\textbf{Perspectiva}     & Terceira pessoa \\
\textbf{Ambientação}     & Vila medieval tomada por mortos-vivos \\
\textbf{Estrutura}       & Mapa único, de 1 a 4 jogadores (co-op online), uma vida por tentativa \\
\textbf{Plataforma}      & PC (Windows e Linux) \\
\textbf{Engine}          & Godot 4.7 (GDScript) \\
\textbf{Público-alvo}    & Jogadores experientes, 16+, acostumados a ação e a escolhas com custo e benefício \\
\bottomrule
\end{tabularx}

% ===============================================================
\section{Conceito}

Jogo de ação 3D em terceira pessoa ambientado em uma vila medieval tomada por mortos-vivos. O jogador escolhe entre quatro aventureiros e recebe de um Lorde a missão de eliminar os inimigos e derrotar o boss responsável pela infestação.

A proposta é permitir que o jogador explore livremente a vila, enfrente grupos de inimigos e tente concluir a missão utilizando apenas uma vida por tentativa.

% ===============================================================
\section{História}

Uma vila medieval foi completamente dominada por mortos-vivos. Diante da ameaça, um Lorde convoca aventureiros para eliminar a infestação e derrotar a criatura responsável pelo ataque.

Caso o aventureiro morra, ele pode aceitar uma nova tentativa da missão. Se decidir desistir definitivamente, o Lorde considera o abandono uma falha grave e ordena sua execução.

\subsection{O contrato}

O texto que o aventureiro assina resume as regras do jogo em cláusulas. O pagamento vem adiantado, contra a assinatura.

\begin{description}[leftmargin=0pt]
    \item[I. Do objeto.] Limpar a vila de todo morto-vivo e derrubar a criatura que os levanta. A vila só é dada por limpa quando o último grupo cai e o Necromante não se levanta mais.
    \item[II. Da vida.] O Contratado entra com uma vida. Caindo, pode aceitar nova tentativa: a vila volta a ser ocupada e cada morto retorna ao seu posto, em outro lugar.
    \item[III. Da bênção.] A cada nova tentativa o Lorde concede uma bênção, à escolha entre duas. A bênção é do contrato, não da pessoa: vale para todos dentro da vila, vivos ou não.
    \item[IV. Da companhia.] Até quatro podem assinar juntos, cada um com seu ofício. Quem cair pode ser erguido por um companheiro que lhe dê a própria poção. Caindo todos, escolhem juntos a bênção; no empate, decide a sorte.
    \item[V. Do abandono.] Desistir é quebra de contrato. A pena é a execução, em praça pública, diante da vila.
\end{description}

As letras miúdas preparam a revelação do calabouço: ``o Lorde não responde pelo que se levantar depois de morto dentro de suas terras, contratados inclusive''.

\subsection{A Companhia da Lanterna}

Cerca de um ano antes do grupo atual, cinco aventureiros aceitaram um contrato do mesmo Lorde para limpar o calabouço sob a vila. A companhia foi fundada por Sir Aldren depois de deixar a ordem de cavalaria, quando ela passou a servir a quem pagava mais. O nome vem da lanterna de ferro que ele levava no cinto: nas tavernas, servia para achar o caminho de volta; na verdade, para achar o espadachim, que se perdia em qualquer corredor com mais de uma porta. O pagamento adiantado deixou Aldren desconfiado desde o primeiro dia (``Lorde nenhum paga adiantado'').

A expedição durou cinco dias e ninguém saiu. Cada um morreu na trava da própria classe, e o que deixaram é a história que o grupo atual encontra no calabouço escondido (em construção):

\begin{description}[leftmargin=0pt]
    \item[Sir Aldren, cavaleiro e fundador.] Metódico, escrevia tudo no diário. Lutava com duas espadas; largou uma só para atravessar o corredor de flechas de escudo erguido e desarmou a armadilha para os outros; morreu no arsenal. Nunca teve coragem de derrubar o espadachim depois que ele voltou morto, e a última página pede isso a quem vier.
    \item[Brakka, bárbaro.] Barulhento e generoso, urrava antes de cada luta e as crianças da vila o seguiam pedindo o urro. Derrubou a parede rachada, achou a adega do Lorde e morreu ali. O esqueletinho da sala da guarda era uma dessas crianças.
    \item[Maelis, maga.] Expulsa da academia por estudar necromancia de perto demais. Entrou na biblioteca pelo portal, a grade desceu atrás dela, e nos livros do Lorde descobriu que quem morre ali levanta e que esse era o plano.
    \item[Wren, ladina.] Batedora de poucas palavras, deduziu os dois selos e a galeria de volta à capela. Errou o último virote na corda da ponte e morreu no abismo.
    \item[O espadachim (Mayoi Michio).] Aldren o achou perdido numa encruzilhada, discutindo com uma placa. Voltou sem pulso, rindo como sempre, e ainda vaga pelos três andares sem lembrar o nome da companhia. Foi ele quem levou a última página de Aldren até a sala secreta, sem saber o que carregava; por isso a página está no andar 3 e os restos de Aldren no arsenal. Em vida tinha uma espada só. Depois que os companheiros caíram, pegou as duas espadas de Aldren, e é daí que vem o estilo de três espadas: uma em cada mão e uma na boca.
\end{description}

\subsection{O grupo atual}

Cada aventureiro jogável tem um elo pessoal com o antecessor da mesma classe, então a trava de cada classe no calabouço é também um reencontro. No jogo, a seleção e o placar continuam mostrando a classe; o nome aparece na história, nos bilhetes e nas falas.

\begin{tabularx}{\textwidth}{@{}lX@{}}
\toprule
\textbf{Aventureiro} & \textbf{Origem, motivo e elo com a Lanterna} \\
\midrule
Ser Edric (Knight) & Escudeiro de Sir Aldren, armado cavaleiro por ele; ficou na ordem quando o mestre saiu e se arrepende. Aceitou o contrato para ir atrás dele. O diário no arsenal é a letra de Aldren, e a última página pede o que Edric jurou nunca fazer: derrubar um companheiro. Arco: do cavaleiro que cumpre qualquer juramento ao que escolhe quais merecem ser cumpridos. \\
Torvald (Barbarian) & Primo mais novo de Brakka, do mesmo clã. Diz que veio pelo pagamento e pela briga; veio também saber por que o primo parou de mandar notícias. O urro dele é igual ao de Brakka, e o esqueletinho reconhece na hora. Arco: do brigão que veio pelo ouro ao que termina o contrato urrando por dois. \\
Ilsa (Mage) & Aluna de Maelis, expulsa da academia um ano depois da mestra pelo mesmo motivo. Seguiu as anotações dela até o Lorde. Atravessa a mesma grade e lê o mesmo aviso. Arco: da estudante fascinada pela necromancia à que entende para que o Lorde a usa. \\
Nessa (Rogue) & Caçadora de recompensas, rival e depois parceira de Wren, que saiu devendo uma aposta. Diz que veio cobrar; veio porque Wren nunca deixava de responder uma carta. Corta a corda que Wren errou (``O seu não vai errar.''). Arco: da mercenária que só confia no que está escrito à que lê as letras miúdas para o grupo. \\
\bottomrule
\end{tabularx}

Para cumprir a cláusula, o grupo inteiro tem que derrubar o espadachim, o companheiro que a Lanterna não teve coragem de derrubar.

% ===============================================================
\section{Gameplay}

O jogador explora livremente a vila e enfrenta grupos de mortos-vivos espalhados pelo cenário.

Cada tentativa começa com o jogador e os grupos de inimigos posicionados aleatoriamente pelo mapa. O jogador possui apenas uma vida.

Ao morrer, pode escolher entre \textbf{Reviver} ou \textbf{Desistir}. Ao reviver, todos os inimigos retornam, novas posições são sorteadas e o jogador recebe uma escolha entre dois buffs aleatórios.

Não existe sistema de níveis para jogadores ou inimigos.

\textbf{Em grupo.} Até quatro jogadores enfrentam a mesma vila, cada um com um aventureiro diferente. Jogar sozinho é um grupo de um: o fluxo é o mesmo. Quem cai espera os aliados, que podem revivê-lo com uma poção guardada. Quando o grupo inteiro cai, todos votam numa das duas bênçãos sorteadas (empate é sorteado) e a tentativa recomeça junta.

\subsection{Ciclo de jogo}

\begin{center}
\begin{tikzpicture}[
    node distance=0.9cm and 1.2cm,
    etapa/.style={draw=Destaque, rounded corners, thick, minimum width=3.4cm,
                  minimum height=0.9cm, align=center, font=\small},
    fim/.style={etapa, fill=Destaque!12},
    seta/.style={-{Stealth[length=2.5mm]}, thick, Destaque},
    rotulo/.style={font=\footnotesize, midway}
]
    \node[etapa] (contrato) {Aceitar o contrato\\(sozinho ou em grupo)};
    \node[etapa, below=of contrato] (sorteio) {Sortear posições do\\jogador e dos inimigos};
    \node[etapa, below=of sorteio] (combate) {Explorar a vila e\\enfrentar os grupos};
    \node[etapa, below=of combate] (boss) {Derrotar o boss};
    \node[fim, below=of boss] (vitoria) {Vitória};

    \node[etapa, right=2.2cm of combate] (morte) {Todo o grupo cai};
    \node[etapa, above=of morte] (buff) {Votar em 1 entre\\2 buffs aleatórios};
    \node[fim, below=of morte] (execucao) {Desistir:\\execução};

    \draw[seta] (contrato) -- (sorteio);
    \draw[seta] (sorteio) -- (combate);
    \draw[seta] (combate) -- (boss);
    \draw[seta] (boss) -- (vitoria);
    \draw[seta] (combate) -- (morte);
    \draw[seta] (boss.east) -- (morte.south west);
    \draw[seta] (morte) -- node[rotulo, right] {Reviver} (buff);
    \draw[seta] (buff) |- (sorteio);
    \draw[seta] (morte) -- (execucao);
\end{tikzpicture}
\end{center}

\subsection{Mecânicas}

\begin{tabularx}{\textwidth}{@{}lX@{}}
\toprule
\textbf{Mecânica} & \textbf{Descrição} \\
\midrule
Movimentação & Terceira pessoa, com exploração livre da vila. Resposta rápida: atinge a velocidade de andar (4\,m/s) em cerca de 0,1\,s e para em cerca de 0,1\,s; correr chega a 7,5\,m/s e gasta estamina. Não há pulo: segurar Espaço corre, e soltar volta a andar. \\
Combate & Corpo a corpo e à distância, conforme o aventureiro escolhido. Clique é o ataque rápido; segurar o botão é o ataque pesado (mais lento e mais forte). Cada golpe rápido tira 1 ponto de vida; o alvo reage ao dano e, ao chegar a zero, morre. \\
Energia & Estamina para todos (correr, escudo, golpes físicos) e mana para o Mage (feitiços). Gastar pausa a recarga por um instante; depois ela volta sozinha. O golpe só sai com o custo inteiro disponível. Detalhes em \ref{sec:energia}. \\
Vida & Pontos inteiros: aventureiros têm 10 e mortos-vivos têm 5. Uma barra de vida flutua sobre cada personagem (verde para aventureiros, vermelha para mortos-vivos). Ao levar dano, um rastro claro mostra o pedaço perdido antes de encolher e um número de dano sobe acima da cabeça. \\
Vida única & O jogador possui apenas uma vida por tentativa. Em grupo, um aliado com uma poção no inventário pode reviver quem caiu: perto do caído aparece ``E Reviver'' sobre ele, e ele volta com metade da vida. A poção só é gasta quando o reviver vale: se outro aliado reviveu antes, ela continua no inventário. Se todos caem, abre a tela de morte. \\
Grupos de inimigos & Mortos-vivos organizados em grupos espalhados pelo mapa. \\
Posições aleatórias & A cada tentativa, o jogador renasce em um ponto sorteado e os grupos de inimigos são reposicionados. \\
Reviver ou Desistir & Escolha apresentada quando o grupo inteiro cai. Desistir (só o host) encerra a missão com a execução do aventureiro. \\
Buffs acumulativos & Ao reviver, o grupo vota num entre dois buffs aleatórios; os buffs se acumulam entre tentativas. \\
Buffs dos inimigos & Os mortos-vivos recebem os mesmos buffs escolhidos pelo jogador nas tentativas anteriores. \\
Boss & O Necromante, um mago esqueleto que surge na praça quando todos os grupos caem: mantém distância, dispara rajadas mágicas e levanta mortos-vivos ao seu redor. \\
Objetos interativos & Baús espalhados pela vila: perto de um, aparece o aviso para abrir; aberto, solta uma poção. Tocar na poção a guarda no inventário (até 4); com o inventário cheio ela fica no chão e é pega assim que sobrar espaço, se o aventureiro ainda estiver em cima dela. Em grupo, dois jogadores na mesma poção não levam uma cada: o host decide quem pegou. Um golpe ou projétil de aventureiro quebra a poção que está no chão, e ela se perde. \\
Inventário rápido & Itens guardados para usar quando quiser: a poção de vida cura 3 ao ser usada (não é gasta com a vida cheia) e serve para reviver um aliado. A roda do mouse para baixo troca o item selecionado; segurar E por 0,6\,s bebe (o gole deixa o aventureiro exposto). \\
Objetos quebráveis & Barris, caixotes e barris de bebida quebram com o ataque pesado (golpe ou projétil), soltando lascas; carroças, cercas e tochas aguentam. Evita ficar encurralado contra o cenário. \\
Troca de habilidade & A roda do mouse para cima troca a habilidade do ataque principal. O Mage alterna entre o Projétil mágico e o Portal; os outros aventureiros têm uma habilidade por enquanto. \\
Portal & Magia do Mage, 20\% de mana por portal. O primeiro uso abre a entrada no chão onde a mira aponta (até 15\,m), o segundo abre a saída; o par dura 12\,s e um par novo fecha o anterior. Aventureiro que entra numa ponta sai na frente da outra (0,5\,s de recarga por pessoa); mortos-vivos não atravessam. Projéteis de todos (setas, projéteis mágicos e os do Necromante) atravessam e saem pela frente do outro portal, na direção que a janela mostra. Uma barra acima do anel esvazia até o portal fechar. O véu do portal mostra o outro lado, como uma janela para a saída, acompanhando o ângulo da câmera de quem olha. Com parede logo à frente da saída, o aventureiro sai no próprio portal. Em grupo, o portal aparece para todos e serve para os aliados. \\
Fogo amigo & Opção do host no lobby, ligada por padrão: golpes e projéteis também acertam aliados. \\
Placar & Abates, assistências, mortes e a última queda de cada jogador, segurando Tab. Um aviso no canto mostra quem abateu quem e como cada um caiu; a tela de morte mostra o histórico de dano daquela vida. \\
Som 3D & Passos, golpes, dor, ossos caindo e feitiços saem da posição de quem os faz; a música muda conforme a cena. \\
Vitória & Limpar todos os grupos (a HUD mostra quantos faltam), derrotar o Necromante e ver a tela ``Contrato cumprido''. \\
\bottomrule
\end{tabularx}

\subsection{Controles}

O jogador escolhe em Configurações $\rightarrow$ Controles se joga com \textbf{teclado e mouse}, com \textbf{controle} (layout Xbox; testado com um 8BitDo Ultimate 2) ou no modo \textbf{automático}, o padrão, em que vale o último dispositivo usado. Com um dispositivo fixo, as ações de jogo respondem só a ele; os menus aceitam os dois sempre.

\begin{tabularx}{\textwidth}{@{}lXX@{}}
\toprule
\textbf{Ação} & \textbf{Teclado e mouse} & \textbf{Controle} \\
\midrule
Mover      & W, A, S, D (relativo à câmera) & Analógico esquerdo \\
Correr     & Segurar Espaço & Clicar no analógico esquerdo; corre até soltar o analógico \\
Câmera     & Movimento do mouse (sobre o ombro direito) & Analógico direito \\
Atacar     & Botão esquerdo: clique é o ataque rápido; segurar 0,35\,s dispara o pesado. Um clique pouco antes da recarga terminar fica guardado & RT (mesmas regras) \\
Habilidade & Botão direito: o Knight ergue o escudo enquanto segura; o Barbarian golpeia com o segundo machado & LT \\
Trocar habilidade & Roda do mouse para cima & Direcional para cima \\
Trocar item & Roda do mouse para baixo & Direcional para baixo \\
Interagir  & E: toque abre baú ou revive aliado caído (com poção no inventário); perto de um baú, segurar não bebe; longe, segurar 0,6\,s bebe o item selecionado. Perto de um aliado caído sem poção, o aviso diz ``Sem poção para reviver'' & X (mesmas regras) \\
Placar     & Tab (segurar) & Back (segurar) \\
Pausar     & Esc & Menu (Start) \\
Menus      & Mouse, setas, Enter, Esc; teclas 1--4 na seleção e 1--2 nas bênçãos & Direcional, A confirma, B volta, LB/RB trocam aventureiro ou aba \\
\bottomrule
\end{tabularx}

As ligações dos dois dispositivos ficam no Input Map do projeto e podem ser trocadas no editor. Sensibilidade do mouse e do analógico e a inversão do eixo vertical também ficam nas configurações. A cena \texttt{scenes/tools/gamepad\_tester.tscn} mostra ao vivo os botões, analógicos, gatilhos e ações recebidos, para conferir um controle novo.

\textbf{Problema conhecido:} um 8BitDo Ultimate 2 ligado por Bluetooth aparece como conectado, mas não manda nenhum comando ao Godot 4.7, embora funcione no painel de controles do Windows. A causa é o driver de 8BitDo da biblioteca SDL, que o Godot usa para ler controles. A solução é definir a variável de ambiente \texttt{SDL\_JOYSTICK\_HIDAPI\_8BITDO=0} (\texttt{setx SDL\_JOYSTICK\_HIDAPI\_8BITDO 0} e reabrir a Steam e o Godot) ou ligar o controle pelo receptor 2.4G ou por cabo.

\subsection{Buffs}

Ao reviver, o grupo vota numa entre duas bênçãos sorteadas (sozinho, o voto é só seu). Pelo contrato do Lorde, a bênção escolhida vale para o aventureiro \emph{e também para todos os mortos-vivos} nas tentativas seguintes. Como o aventureiro tem 10 de vida e enfrenta vários mortos-vivos de 5 de vida, o mesmo buff pesa diferente para cada lado: a escolha é decidir qual efeito favorece mais o jogador do que os inimigos. Metade dos buffs tem custo e benefício, pensados para o público experiente.

\begin{tabularx}{\textwidth}{@{}lX@{}}
\toprule
\textbf{Buff} & \textbf{Efeito (em todos os personagens)} \\
\midrule
Lâmina Afiada & +1 de dano \\
Couraça & +3 de vida máxima \\
Passos Leves & +20\% de velocidade \\
Frenesi & $-$25\% de tempo entre ataques, $-$2 de vida máxima \\
Fúria Cega & +1 de dano, $-$3 de vida máxima \\
Sangue Frio & +2 de vida máxima, $-$10\% de velocidade \\
\bottomrule
\end{tabularx}

Regras: cada buff acumula até 3 vezes; os dois sorteados são sempre diferentes; vida máxima nunca fica abaixo de 1; velocidade e tempo entre ataques nunca caem abaixo de 20\% do valor base. Desistir na tela de morte, ou abandonar o contrato pela pausa, leva à cena de execução do aventureiro e depois ao menu principal; o próximo contrato começa do zero (tentativas e buffs zerados).

\subsection{Energia: estamina e mana}
\label{sec:energia}

Todo aventureiro tem uma barra de estamina de 100; o Mage tem também uma de mana. Gastar pausa a recarga (0,9\,s na estamina, 1\,s na mana); depois a barra volta sozinha. Se a estamina zerar, o aventureiro fica exausto: não ataca, não usa o escudo e não corre até a barra encher de novo por completo, e a barra fica apagada na HUD enquanto isso. A mana do Mage não trava ao zerar. A recarga parte de uma velocidade base comum a todos (24 por segundo, baixada de 30 depois de medir que a barra cheia voltava rápido demais para o ataque contínuo) e cada classe tem um multiplicador próprio. Correr ou segurar o escudo gasta estamina continuamente e impede a recarga enquanto dura. Barbarian, Mage e Rogue, quando zeram a estamina ou a mana, já começam a recarregar na hora, sem a pausa; o Knight mantém a pausa. Quem zera a estamina correndo volta a andar e só corre de novo com a barra em 30\%.

\begin{tabularx}{\textwidth}{@{}lccccX@{}}
\toprule
\textbf{Aventureiro} & \textbf{Paga com} & \textbf{Rápido} & \textbf{Pesado} & \textbf{Recarga} & \textbf{Outros gastos} \\
\midrule
Knight    & estamina & 15 & 35 & 1$\times$ & escudo: 6 por segundo e 12 por golpe defendido \\
Barbarian & estamina & 15 & 50 & 1$\times$ & --- \\
Rogue     & estamina & 12 & 30 & 1$\times$ & --- \\
Mage      & mana     & 20 & 45 & mana 1,5$\times$, estamina 0,5$\times$ & --- \\
\bottomrule
\end{tabularx}

Todos gastam 12 de estamina por segundo correndo. Com menos energia que o custo, o golpe não sai e a barra pisca em vermelho. Os números ficam no inspetor de cada aventureiro (grupos ``Energia'' e ``Pesado'').

% ===============================================================
\section{Personagens}

O jogador poderá escolher entre quatro aventureiros, com diferentes estilos de combate, incluindo personagens corpo a corpo e de longo alcance.

\begin{tabularx}{\textwidth}{@{}lX@{}}
\toprule
\textbf{Personagem} & \textbf{Papel} \\
\midrule
Aventureiros & Personagens controlados pelos jogadores. \\
Lorde & Responsável por entregar a missão. \\
Mortos-vivos & Inimigos espalhados pela vila. \\
Boss & Principal inimigo e responsável pela infestação. \\
\bottomrule
\end{tabularx}

\subsection{Aventureiros jogáveis}

Cada aventureiro tem um ataque único. Todos atacam na direção em que a câmera olha. Mage e Rogue têm uma mira no centro da tela: o projétil sai da arma em direção ao ponto que a mira cobre (parede, chão ou morto-vivo, até 40\,m), inclusive para cima ou para baixo, então acertar de longe exige apontar a câmera; o ataque rápido causa 1 de dano por acerto (mais os buffs) e só acerta outro aventureiro com o fogo amigo ligado. Projéteis param ao bater em paredes e deixam rastro: o projétil do Mage queima uma marca violeta na superfície (some após 6\,s) e o virote do Rogue fica cravado onde bateu, na parede ou no corpo do morto-vivo, acompanhando-o enquanto ele se move (some após 4\,s). Golpes corpo a corpo não atravessam paredes: só acertam quem está com o caminho livre até o atacante. O virote da besta fica cravado no osso mais próximo de onde acertou e acompanha o morto-vivo andando, girando e caindo ao morrer.

Regras comuns de ataque, pensadas para o controle não ficar pesado:
\begin{itemize}
    \item durante o ataque o aventureiro ainda se move a 30\% da velocidade, sem girar, para manter a mira;
    \item logo depois do impacto, mover ou atacar de novo corta o resto da animação; parado, a animação termina normalmente;
    \item ser atingido durante a preparação (antes do impacto) cancela o golpe; depois do impacto, o golpe já contou;
    \item ao ser atingido, o aventureiro recupera o controle em 0,3\,s;
    \item o morto-vivo que cambaleia fica 1,4\,s sem cambalear de novo: ele continua levando dano e sendo empurrado, mas termina o próprio ataque, então bater sem parar não o deixa preso (medido: antes, um esqueleto apanhando sem parar nunca acertava o aventureiro corpo a corpo).
\end{itemize}

Peso do golpe: quando um golpe acerta, a animação de quem bateu e de quem apanhou para por um instante (0,06\,s no rápido, 0,11\,s no pesado), a câmera de quem bateu treme de leve (mais no pesado) e toca um som de confirmação só para ele, mais grave no pesado. Quem apanha é empurrado para longe de quem bateu: o golpe rápido empurra pouco, para não tirar o alvo do alcance do combo, e o pesado empurra o dobro; o minion recua bastante, o elite e os aventureiros menos, e o Necromante quase nada. Golpe bloqueado pelo escudo não empurra, mas quem bateu ainda sente o tranco (pausa, tremor e som). A pausa e o tremor são só visuais; em rede, a pausa em quem apanhou numa máquina remota é quase imperceptível, porque a animação de dor chega pela rede logo depois; o empurrão é calculado por quem controla o personagem e chega aos outros pela rede.


\begin{tabularx}{\textwidth}{@{}llX@{}}
\toprule
\textbf{Aventureiro} & \textbf{Alcance} & \textbf{Estilo de combate} \\
\midrule
Knight & Corpo a corpo & Espada e escudo de torre. \textbf{Combo de três golpes} (vertical, diagonal, horizontal) que segue enquanto o jogador continua atacando e reinicia após uma pausa; acerta só à frente; recarga de 0,6\,s. \textbf{Pesado}: golpe de cima com as duas mãos, dano dobrado. \textbf{Escudo} (botão direito): enquanto segura, bloqueia golpes e projéteis que chegam pela frente (arco de 140°), fica parado (gira com a câmera) e não ataca; gasta estamina e, sem estamina, a guarda quebra; o bloqueio aparece como ``Bloqueado'' sobre a barra de vida \\
Barbarian & Corpo a corpo & Dois machados. \textbf{Machado}: uma machadada só à frente; recarga de 0,6\,s. \textbf{Pesado (giro)}: segurando o ataque, gira com os dois machados e acerta todos os mortos-vivos em volta, com dano dobrado; custa metade da estamina e tem recarga de 1,6\,s. Continuando a segurar, ele emenda giros curtos (20\% da estamina por volta) e anda a 70\% da velocidade enquanto gira; soltar o botão, ficar sem estamina ou apanhar encerra o giro (depois de apanhar, é preciso soltar e segurar de novo). \textbf{Segundo machado} (botão direito): golpe rápido só à frente, recarga de 0,5\,s, para alternar com o machado principal \\
Mage & Longo alcance & Cajado, usa mana. \textbf{Projétil mágico}: mais lento (14\,m/s), alcance de 6,91\,m e \emph{atravessa} os inimigos em linha; recarga de 1,0\,s. \textbf{Pesado}: feitiço longo, 2,5$\times$ o dano. \textbf{Portal} (troca na roda do mouse): abre um par de portais para o grupo atravessar \\
Rogue & Longo alcance & Besta. \textbf{Virote}: rápido (30\,m/s), alcance de 6\,m (um pouco menos que o Mage), para no primeiro inimigo; recarga de 0,7\,s. É o mais rápido: anda a 4,8\,m/s e corre a 9\,m/s, 20\% acima dos outros. \textbf{Pesado}: disparo mais lento, dano dobrado \\
\bottomrule
\end{tabularx}

\subsection{Mortos-vivos e IA}

Os mortos-vivos não se limitam a andar e bater: eles percebem, reagem com atraso humano e, nas próximas fases, se coordenam em grupo e respondem às ações do jogador. Todos usam a mesma base de física e animação do jogador e são mais lentos que ele (andar 2,8\,m/s, correr 5\,m/s), para que fugir seja uma opção real.

\textbf{Percepção.} Em grupo, cada morto-vivo mira o aventureiro vivo mais próximo. Visão em cone de 120\textdegree{} até 14\,m, bloqueada por paredes; audição de corrida e de ataques num raio de 6\,m, mesmo sem ver; percepção de proximidade: nota o jogador a até 2,5\,m mesmo de costas, se não houver parede; ser atingido, inclusive por um disparo de longe, revela a posição do jogador; memória da última posição do jogador por 3\,s.

\textbf{Comportamento (implementado).}
\begin{itemize}
    \item \textbf{Adormecido}: deitado no chão; não enxerga, mas desperta com o jogador a 7\,m, com barulho ou ao ser atingido. Enquanto dorme não mostra barra de vida; ela aparece quando a animação de despertar avisa que o perigo começou. Atingido deitado, começa a levantar na hora. Deitado ou levantando, leva o dano normalmente mas não se encolhe com o golpe, então a animação de levantar nunca é cortada.
    \item \textbf{Em guarda}: ao perceber o jogador, espera 0,3\,s antes de agir, para não parecer onisciente.
    \item \textbf{Perseguição}: segue o jogador pela navegação, contornando paredes; corre quando está longe.
    \item \textbf{Ataque}: só golpeia quase de frente para o jogador, com aviso de 0,6\,s entre o início do golpe e o impacto (tempo para esquivar ou cortar o golpe); recarga de 1,35\,s.
    \item \textbf{Busca}: ao perder o jogador, vai até a última posição conhecida, procura por 2,5\,s e volta à guarda.
    \item \textbf{Corpo}: o morto-vivo derrubado fica caído por 6\,s, afunda no chão e sai da cena, deixando só uma mancha escura onde caiu. Mantém a vila limpa e reduz o que precisa ser desenhado.
    \item \textbf{Elite}: mortos-vivos podem formar um grupo. Quando o último do grupo cai, os corpos ficam parados por 2\,s, deslizam até o ponto do grupo, afundam juntos e dali se levanta um esqueleto elite: maior, com elmo de chifres e aura roxa, 15 de vida, golpe de 2 de dano, visão e audição maiores. Ele já levanta sabendo onde o jogador está e vai atrás dele. Limpar uma área tem um último confronto, e o elite serve de base para o boss. Para montar um grupo na fase, basta colocar os mortos-vivos como filhos de um nó \texttt{EnemyGroup}: cada um só pode ter um pai, então grupos não se misturam. O elite nasce no ponto onde o nó do grupo foi colocado; só os corpos a até 8\,m desse ponto são puxados (quem caiu longe some como um corpo comum), e o editor mostra um aviso se algum membro for posicionado longe demais.
\end{itemize}

\textbf{Boss: o Necromante.} Um esqueleto mago maior, com cajado e aura verde, 30 de vida. Surge na praça, levantando do chão, quando o último grupo da vila cai. Usa o mesmo cérebro dos outros mortos-vivos, mas com alcance de 14\,m, mais que o dos aventureiros a distância: fica longe e dispara rajadas mágicas (2 de dano, recarga de 2\,s). A cada 9\,s de caça levanta dois esqueletos ao seu redor, com no máximo quatro vivos; quando ele morre, os invocados desabam junto. A HUD mostra uma barra de vida grande com o nome dele.

\textbf{Próximas fases.}
\begin{enumerate}
    \item \textbf{Grupo}: no máximo alguns mortos-vivos atacam ao mesmo tempo, os demais cercam o jogador por lados diferentes; quem vê o jogador alerta os vizinhos.
    \item \textbf{Reações}: ler o golpe do jogador; bloqueio frontal (Warrior com escudo); esquiva; atacar quando o jogador erra um golpe. As chances de reação crescem 5\% por tentativa, até 60\%.
    \item \textbf{Variantes}: Warrior (escudo), Rogue (besta que mira prevendo o movimento e mantém distância) e Mage (projétil), cada um com seu perfil de números.
\end{enumerate}

% ===============================================================
\section{Mundo}

O jogo acontece em uma vila medieval explorável, composta por ruas, casas, praça, construções e áreas abertas.

A missão principal acontece em um único mapa, a vila. Debaixo dela fica o calabouço escondido, uma área bônus que ainda será ligada à vila por uma passagem secreta.

A cada nova tentativa, tanto o jogador quanto os grupos de inimigos recebem novas posições aleatórias dentro da vila, fazendo com que os encontros e o caminho percorrido possam mudar entre as partidas.

\textbf{Pedraalta, a vila da montanha (fase atual).} Uma vila mineira de prata na encosta: o Lorde quer a mina de volta, não a vila, e por isso contrata aventureiros descartáveis. O mapa tem 44\,m $\times$ 122\,m e sobe de sul para norte em terraços (cotas 0, +3, +6, +9 e +12\,m), ligados por rampas de 4\,m de largura e até 25 graus que alternam de lado para obrigar a cruzar cada terraço. As bordas são visíveis: paredão de rocha a oeste e ao norte, portão desmoronado ao sul e cerca sobre o precipício a leste. São cinco zonas com um grupo de mortos-vivos cada, sorteado entre os pontos da zona; a zona da entrada sempre recebe o grupo dos dormentes, o mais fácil, para o primeiro encontro ensinar: \textbf{A, Porteira do Cais de Carga} (portão arrombado por dentro, carroças viradas, bandeiras do Lorde); \textbf{B, Terraços dos Fornos} (casas, forno frio, varal com roupa e camas na rua: dormiram e não acordaram; fonte como respiro); \textbf{C, Pedreira do Corte Fundo} (ferramentas largadas, entulho, um veio de minério claro demais no paredão e a entrada gradeada do calabouço, com a lanterna da Companhia quebrada na boca); \textbf{D, Capela dos Mil Degraus} (covas abertas de dentro para fora e velas ainda acesas; a capela não é refúgio); \textbf{E, Varanda do Vigia} (ramal opcional pela ponte sobre a fenda, com o posto do vigia e um baú). No topo, \textbf{F, o Mirante do Último Sino}: a arena do Necromante, um círculo de pedra vermelha com seis braseiros, o sino rachado e o cadafalso do Lorde. Há um baú em B, C, E e F. Sem seta de objetivo, os marcos guiam: a fumaça e o brilho da pedreira, as velas da capela e os braseiros vermelhos do topo. A vila plana anterior continua no projeto (\texttt{village\_level.tscn}) como referência.

\textbf{A vila (primeira versão).} Um vilarejo medieval de 48\,m $\times$ 48\,m, à noite, iluminado por tochas nas paredes. Todo o chão é de pedra (tijolo, tijolo vermelho e calçamento irregular). No centro, uma praça calçada com carroças, caixotes e barris que servem de cobertura; em volta, oito casas de pedra e enxaimel, de um ou dois andares, formando ruas e becos; ao sul, um campo aberto cercado. As portas em arco são vãos de verdade: dá para entrar nas casas, que viram esconderijos e atalhos.

\textbf{Sorteio dos encontros.} O mapa tem seis pontos de nascimento do jogador e treze pontos de grupo. A cada tentativa o jogador nasce num ponto sorteado e cinco grupos ocupam pontos a pelo menos 16\,m dele. Os grupos são cenas prontas: \textbf{patrulha} (2 esqueletos, sem elite), \textbf{guarda} (3 esqueletos, com elite) e \textbf{dormentes} (3 esqueletos deitados, com elite). Um teste garante que de todo ponto de grupo se anda até todo ponto de nascimento. Em grupo, todos usam a mesma semente de sorteio, então a vila sai igual em todas as máquinas, e os jogadores nascem lado a lado no ponto sorteado.

\subsection{O calabouço escondido (em construção)}
\label{sec:calabouco}

Masmorra escura de três andares (cerca de 90\,m $\times$ 120\,m), montada com as peças do KayKit Dungeon em grade de 4\,m, pensada para uns 10 minutos: descer, abrir a descida, chegar ao tesouro e voltar à saída. Sem céu: fundo preto, luz ambiente baixa, neblina volumétrica e tochas que tremulam são a única luz; o nicho do espadachim fica sem tocha nenhuma. \textbf{Andar 1}: entrada, salão de pilares, nicho escuro, sala da guarda (dois esqueletos jogando dados), a biblioteca atrás de uma grade e a adega atrás de uma parede rachada. \textbf{Andar 2}: cripta, ossuário (dormentes), corredor das lâminas, sala dos espinhos, capela (guardas) e o corredor de flechas até o arsenal. A descida para o andar 3 fica atrás do \textbf{portão dos dois selos}: uma alavanca no ossuário e outra na capela, as duas guardadas por mortos-vivos. \textbf{Andar 3}: corredor dos dardos, antessala, subida do barril, sala do tesouro (com o mímico), o abismo da ponte levadiça e a sala secreta com o portão selado (``Além deste portão, outro calabouço espera.''). O chão do calabouço tem colisão de caixa lisa no topo de cada ladrilho: a malha exata, com sulcos e pedrinhas, prendia o personagem nas bordas. Seis grupos de mortos-vivos ocupam os andares, com malha de navegação própria, feita com portões e paredes abertos para que quem persegue atravesse o que o grupo já abriu.

\textbf{Objetos interativos (Area3D).}
\begin{itemize}
    \item \textbf{Poção de mana}: recupera 40 de mana; só quem usa mana (o Mage) consegue pegar, os outros a deixam no chão.
    \item \textbf{Espinhos}: sobem e descem num ciclo de 4\,s; cada subida fere uma vez quem está em cima, inclusive mortos-vivos, e o escudo não protege de baixo. As armadilhas vizinhas começam defasadas para formar um ritmo de travessia.
    \item \textbf{Lâmina pendular}: balança de um lado a outro (3 de dano); no fim do balanço deixa o meio livre.
    \item \textbf{Dardos}: pisar na placa dispara três virotes da parede (recarga de 1,5\,s); atrair um morto-vivo para a placa também o fere.
    \item \textbf{Barril rolando}: inspirado no level design de Dark Souls. Subir a escada solta um barril lá de cima, um corpo rígido com impulso e torque iniciais que desce rolando pela física e atropela quem estiver no caminho (2 de dano e empurrão), recarga de 4\,s.
    \item \textbf{Saída}: segurar E na saída encerra o contrato (ver abaixo).
    \item \textbf{Alavanca, bilhete e saída}: todos mostram ``E <ação>'' para quem chega perto. Alavanca e saída passam pelo host, que confere o alcance; o bilhete abre só na tela de quem leu e fecha ao se afastar.
\end{itemize}

\textbf{Travas por classe.} O caminho principal e todos os mortos-vivos da cláusula se fazem com qualquer aventureiro; cada classe abre para o grupo um atalho ou uma sala com loot e história, para que em grupo cada um tenha seu momento.
\begin{itemize}
    \item \textbf{Mage, a grade da biblioteca}: a grade segura o corpo, mas deixa passar a mira, os tiros e o portal (fica numa camada própria, \texttt{grades}). O Mage abre a saída do portal do lado de dentro e puxa a alavanca que ergue a grade para todos.
    \item \textbf{Knight, o corredor de flechas}: como no Dark Souls, o chão do corredor tem placas de pressão; pisar numa estala e, 0,5\,s depois, as seteiras da parede do fundo do arsenal disparam pelo corredor. Como vêm de frente, o escudo apara cada flecha (gasta estamina): o Knight pisa, ergue o escudo, avança. Na parede leste do arsenal, a alavanca desarma o corredor.
    \item \textbf{Rogue, a corda da ponte}: no abismo, a ponte levadiça erguida fecha a passagem e há um fosso de lâminas na frente. A corda que a segura fica a mais de 4\,m de altura: só o virote da besta corta (a magia não corta e a lâmina não alcança). A ponte desce sobre o fosso e abre a galeria, que sobe de volta à capela: um atalho para a volta, no estilo de Dark Souls.
    \item \textbf{Barbarian, as paredes rachadas}: só o giro pesado do Barbarian as derruba (marcado no golpe, \texttt{breaks\_walls}). Uma esconde a adega do andar 1; a outra, a sala secreta do andar 3.
\end{itemize}

\textbf{A Companhia da Lanterna.} A história vem pelo cenário: dez bilhetes e os restos de quatro aventureiros (cavaleiro, bárbaro, maga e ladina) que pegaram o mesmo contrato antes do grupo, cada um morto na trava da própria classe. A ordem do Lorde na entrada repete a cláusula; o diário de Sir Aldren conta que eram cinco, contando o espadachim que vivia se perdendo; Maelis descobre que quem morre ali levanta e que o Lorde sabia; o ossuário está cheio de crânios com marca de contrato; e a última página, na sala secreta, pede o que a cláusula exige: derrubar também o espadachim. Ele mesmo, no andar 3, se lembra do nome da companhia.
Toda armadilha nomeia a própria causa no histórico de dano. Em grupo, cada máquina fere só quem controla, e armadilha sem dono é decidida pelo host.

\textbf{A cláusula do contrato.} O Lorde mandou limpar tudo, e a HUD mostra ``Mortos-vivos restantes''. Sair com zero cumpre o contrato e abre a vitória; sair com algum de pé é falha, e o aventureiro vai à execução. Quando o calabouço for ligado à vila, a mesma conta passa para a vitória contra o Necromante.

\textbf{O espadachim perdido.} Easter egg inspirado em Zoro (One Piece) e no carisma de Solaire (Dark Souls): um esqueleto de bandana verde com três espadas na cintura, perdido. O grupo o encontra no escuro; quem atira nele ouve ``Uepa! Pra que essa agressividade toda? Vamos com calma.'', e ele sai da sombra sem levar dano. Ele se chama \textbf{Mayoi Michio} (sobrenome primeiro, como no japonês): Michio é ``homem do caminho'' e Mayoi, ``se perder''. A primeira conversa, em qualquer andar, é a apresentação; até lá ele é só o ``Espadachim perdido'' no histórico de dano e nos avisos, e depois dela passa a ser Mayoi Michio. Perto dele aparece ``E Conversar''; cada aperto traz a próxima fala, que fica na tela até a seguinte (o resto do combo depois do susto não corta a fala: ele tem 2\,s de folga). A conversa acompanha o grupo: primeiro o cumprimento do andar, depois o assunto mais importante que já vale, e por fim as falas do andar, em rodízio. Os assuntos ficam na cena da fase (\texttt{DialogueTopic}): cada um observa uma propriedade do mundo (portão aberto, selos puxados, corda cortada, parede derrubada, mímico morto, tesouro aberto, mortos-vivos restantes) e só aparece dentro da faixa marcada. Na ordem: \textbf{lembranças} (abrir a sala de um antigo companheiro faz ele lembrar um nome solto: ``Maelis! Uma moça que brigava com os livros e ganhava''), \textbf{o que o grupo fez} (``Um selo já foi! Falta o outro'') e \textbf{indiretas sobre o que falta} (a biblioteca trancada, a parede que cheira a cerveja, a corda alta demais para a espada). As falas se ajustam ao grupo: sozinho, ele fala com você; em grupo, com vocês. \textbf{Mais um no grupo.} Depois da primeira rodada de falas do andar, ele pergunta ``O que você acha de mais um no grupo?''. Perto dele aparecem E Aceitar e F Recusar (X e B no controle); recusar faz ele perguntar de novo no andar seguinte. Aceito, sozinho ou em grupo, ele segue quem disse sim (ou o aventureiro mais perto, se esse cair), saca uma espada e corta o morto-vivo mais próximo. Ele também atrapalha: com o fogo amigo ligado, as espadas dele pegam o grupo; de vez em quando jura que o caminho é por outro lado (``É por aqui! Confia em mim.''), some até alguém ir buscá-lo ou falar com ele; e pisa em toda placa do caminho, disparando armadilha em cima do grupo (armadilha não fere ele, só a reclamação). Se o grupo sobe de andar com ele perdido longe, ele aparece na frente, como sempre. Continua contando na cláusula: apanhar do grupo só magoa (``Ei! Eu tô do seu lado!''), mas com a vida baixa ele sai do grupo e a luta dos três estilos começa.

A cláusula vira piada que dói: ``Faltam 3 mortos-vivos. Engraçado: toda vez que eu conto, sobra um que eu não acho''. A cada descida ele reaparece, perdido, no andar seguinte. Ele também é um morto-vivo e conta na cláusula: para vencer é preciso derrubá-lo. Com a vida alta ele só reclama; com 3 de vida (de 12) saca uma espada e luta (estilo de uma espada), com 2 saca a segunda (duas espadas) e com 1 a terceira, mordida pelo cabo (Santoryuu, estilo de três espadas). Brigando, cada golpe que ele recebe vale só 25\%, o que o torna um semi-boss; cada estilo tem golpe, recarga e animações próprios, ajustáveis no inspetor.

\textbf{O esqueletinho.} Easter egg na sala da guarda, perto da parede rachada da adega: um esqueleto pequeno que segue quem chega pedindo ``Faz o urro! Faz o urro!''. Só quem urra (o Barbarian) vê ``E Urrar''; o urro faz ele desmontar de emoção, se remontar e pedir de novo, e o primeiro urro dá uma poção de presente e a dica da adega (``O tio grandão que morava atrás daquela parede rachada também urrava assim!''). Para as outras classes ele só pergunta pelo grandão. Apanhar só o desmonta, e ele não conta na cláusula, a não ser que um Barbarian por perto negue o urro seis vezes: emburrado, ele vira um morto-vivo de verdade, entra na conta e pode cair. Um urro tardio faz as pazes e o tira da conta de novo.

\textbf{O mímico.} Na sala do tesouro, ao lado do baú de verdade, há um baú igual que é um mímico. A única pista é a tampa, que de vez em quando abre uma fresta quase imperceptível. Quem tenta abrir ou bate nele é mordido: ele mostra os dentes e a língua e persegue o aventureiro mais perto aos pulos (14 de vida, mordida de 2). Derrotado, solta uma poção de vida. Não é morto-vivo, então não conta na cláusula.

% Para incluir um mapa ou esboço, salve a imagem em docs/img/ e ajuste o nome.
\IfFileExists{img/mapa.png}{%
\begin{figure}[h]
    \centering
    \includegraphics[width=0.8\textwidth]{img/mapa.png}
    \caption{Esboço da vila}
\end{figure}}{}

% ===============================================================
\section{Interface}

Todas as telas compartilham um tema único de madeira escura e dourado. Ações destrutivas usam borda vermelha e sempre pedem confirmação; a ação principal de cada tela vem destacada e recebe o foco do teclado. A interface foi desenhada para 1152$\times$648 e escala com o tamanho da janela, mantendo a proporção.

\textbf{Fluxo de telas:} Menu principal $\rightarrow$ Seleção de aventureiro e Assinatura do contrato (sozinho) ou Contrato em grupo (lobby) $\rightarrow$ Entrada na missão $\rightarrow$ Partida (HUD, Pausa) $\rightarrow$ Morte $\rightarrow$ Reviver (nova tentativa na mesma partida) ou Desistir $\rightarrow$ Execução $\rightarrow$ Menu principal. Abandonar o contrato pela pausa também leva à Execução.

\begin{tabularx}{\textwidth}{@{}lX@{}}
\toprule
\textbf{Tela} & \textbf{Conteúdo} \\
\midrule
Menu principal & Título, ``Aceitar contrato'', ``Contrato em grupo'', Configurações, Sair e créditos, sobre uma cena 3D da vila com câmera girando \\
Contrato em grupo & Nome do jogador, Hospedar ou Entrar (pelo endereço do host). Na sala, em duas colunas: à esquerda os endereços do host (clicar copia) e o grupo; à direita a escolha do aventureiro (sem repetir), o fogo amigo (só o host) e Iniciar (só o host, quando todos escolheram) \\
Seleção de aventureiro & Os quatro aventureiros lado a lado em 3D; o escolhido dá um passo à frente, ganha um anel dourado no chão e o nome destacado. Ficha com estilo, ataque e descrição. Escolha por clique, teclas 1 a 4 ou setas; a escolha vale para todas as tentativas do contrato \\
Assinatura do contrato & Ao aceitar o aventureiro, o contrato do Lorde abre em pergaminho: as cinco cláusulas resumidas, as letras miúdas, o ofício escolhido e a linha ``Assine aqui''. O nome digitado aparece na hora como assinatura inclinada sobre a linha do Contratado; Assinar (ou Enter) começa a missão, Voltar (ou Esc) retorna à escolha. O nome assinado passa pela mesma limpeza do nome do lobby (até 16 caracteres, sem caracteres invisíveis; em branco vira ``Aventureiro''), fica salvo para o próximo contrato e aparece nos avisos, no placar e na vitória (``Contrato de Fulano, cumprido na tentativa 2''; em grupo, todos os nomes). Quem já assinou antes encontra o nome preenchido e só confirma, o que serve também para quem joga no controle \\
Configurações & Áudio (geral, música, efeitos), vídeo (tela cheia, VSync) e controles (sensibilidade do mouse, inverter eixo vertical, lista de teclas). Salvas automaticamente \\
HUD & Segue o estudo de HUD do Figma. No canto superior direito, o objetivo: ``Grupos restantes'' com um losango por grupo, aceso enquanto o grupo vive, e depois ``Derrote o Necromante''. Ele não aponta direção, porque achar os grupos faz parte de explorar a vila. Na primeira tentativa, um aviso no centro lembra que Esc (ou Start) abre a pausa com os controles, e o aviso de que o Necromante despertou fica 5\,s na tela, com som. No canto superior esquerdo, a tentativa em numeral romano com as bênçãos ativas em etiquetas ao lado (nome e acúmulos; o efeito de cada uma fica na pausa), a vida em barra vermelha com número e um rastro claro que mostra por um instante quanto o golpe tirou, e a estamina (âmbar) e a mana (azul, só o Mage), divididas em blocos do custo do golpe atual (bloco que não encheu não pode ser gasto); aviso de abates e quedas no canto e placar segurando Tab. No canto inferior esquerdo, uma cruz de slots com ícones: em cima a habilidade ativa (com o custo em \% num selo âmbar no canto e só o nome ao lado; o resto fica na tela de controles e habilidades da pausa), à esquerda o botão direito (escudo ou segundo machado; some para quem não tem, mantendo o lugar na cruz) e embaixo o item rápido (com a quantidade, o nome e o lembrete ``segure E''; sem item, o slot fica vazio, sem texto). Segurar o ataque ou o E enche um anel dourado em volta do slot correspondente, mostrando quanto falta para o golpe pesado ou o gole. Os ícones são renderizados dos próprios modelos 3D do jogo. Para Mage e Rogue, e com o Portal ativo, uma mira no centro da tela indica o alvo. \\
Barras 3D & Sobre os mortos-vivos (vermelha, com rastro do dano); escondida enquanto o morto-vivo dorme, aparece quando ele começa a levantar. Sobre os aliados, nome e barra verde; sobre um aliado caído, o aviso de reviver para quem tem poção por perto. O jogador usa a barra do HUD. Números de dano sobem de todos os personagens \\
Entrada na missão & Cena curta à noite: o aventureiro escolhido chega pela estrada até a entrada da vila (``Contrato aceito.'', ``O Lorde quer a vila livre dos mortos.''); em grupo, os aliados vêm em formação atrás. Qualquer tecla marca ``pronto''; a partida começa quando todos estão prontos \\
Pausa & Em duas colunas e com pouco texto. À esquerda, o selo, ``Contrato pausado'' e os botões Continuar, Controles e habilidades, Configurações e Abandonar contrato (com confirmação; leva à execução). À direita, só a situação do contrato: tentativa, aventureiro, objetivo atual e as bênçãos ativas com o efeito de cada uma (``valem também para os mortos-vivos''). O botão Controles e habilidades abre uma tela à parte com a tabela de controles, que troca sozinha entre teclado e controle, e a ficha da classe: estilo, descrição e cada habilidade com o custo em \% da barra (rápido e pesado) e o botão direito; Esc volta à pausa. Em grupo, o jogo não congela, mas o aventureiro para de obedecer enquanto o menu está aberto (solta o escudo e cancela o que estava segurado); abandonar tira só quem saiu \\
Morte & Histórico de dano da vida (de quem veio cada golpe, curas e o golpe final), tentativa, bênçãos acumuladas, Reviver ou Desistir (só o host; com confirmação; leva à execução). Reviver mostra duas cartas de bênção para votar, com o mouse ou as teclas 1 e 2 \\
Execução & Cena em duas tomadas: o aventureiro escolhido vai embora pela estrada da vila ao entardecer (``Ele deu as costas ao contrato''); corta para a praça pública, onde ele está ajoelhado num cadafalso diante da multidão, ouve o veredito do Lorde e recebe o golpe de machado do carrasco encapuzado, enquanto a multidão comemora; termina em ``Contrato encerrado'' e volta ao menu. Qualquer tecla pula \\
\bottomrule
\end{tabularx}

\subsection{Áudio}

Todo som de personagem é posicional (3D): passos contados pela distância andada no chão, golpe saindo no momento do impacto, dor ao levar dano, estalo no escudo ao bloquear, ossos caindo na morte do esqueleto, rangido ao levantar do chão e um sino grave quando o elite se forma. Virotes e feitiços fazem barulho onde batem. Cada grupo de sons sorteia entre variações e muda levemente o tom, para não repetir. Os botões de todos os menus têm som de clique e de foco. A música segue a cena: uma melodia medieval no menu e na seleção, um ambiente sombrio na vila e nas cutscenes. Música e efeitos têm volumes próprios nas configurações.

% ===============================================================
\section{Arquitetura técnica}

\subsection{Personagens}

Todos os personagens compartilham uma classe base. Cada classe filha só decide a \emph{intenção} de movimento; física, vida e animação ficam na base.

\begin{center}
\begin{tikzpicture}[
    node distance=0.8cm and 0.6cm,
    classe/.style={draw=Destaque, rounded corners, thick, align=center, font=\small, inner sep=5pt},
    seta/.style={-{Stealth[length=2.5mm]}, thick, Destaque}
]
    \node[classe] (character) {\textbf{Character}\\ física, vida, dano, morte, animação};
    \node[classe, below left=of character] (adventurer) {\textbf{Adventurer}\\ entrada, câmera, energia,\\ máquina de estados};
    \node[classe, below right=of character] (enemy) {\textbf{Enemy}\\ percepção, cérebro, navegação};
    \node[classe, below=of enemy] (boss) {\textbf{Necromante}\\ Enemy + Summoner};
    \draw[seta] (adventurer) -- (character);
    \draw[seta] (enemy) -- (character);
    \draw[seta] (boss) -- (enemy);
\end{tikzpicture}
\end{center}

Cada aventureiro e cada tipo de morto-vivo é uma cena herdada da base, que troca apenas o modelo, as armas visíveis e os atributos. A IA dos mortos-vivos é dividida em percepção (o que sabem), cérebro com máquina de estados (o que decidem, a cada 0,15\,s) e navegação (como se movem); os números de cada tipo ficam num perfil de dados. O ataque é um componente separado, reaproveitado por aventureiros e inimigos: \texttt{Attack} (recarga, animação e momento do impacto) com as variantes \texttt{MeleeAttack} (o formato da hitbox define golpe frontal ou giro) e \texttt{RangedAttack} (dispara um \texttt{Projectile}). As telas apenas emitem sinais; o \texttt{GameSession} decide trocas de cena e reinícios. Buffs são recursos de dados (\texttt{Buff} com uma lista de \texttt{StatModifier}); o estado da missão (tentativa e buffs acumulados) fica no \texttt{RunState}, que sobrevive à recarga da cena a cada Reviver. Quem pode ser atingido é definido pelas camadas de física: mundo, aventureiros e inimigos.

\subsection{Máquina de estados do jogador}

O aventureiro tem um nó \texttt{StateMachine} com um nó filho por estado, no mesmo formato do cérebro dos mortos-vivos: \textbf{Idle}, \textbf{Walk}, \textbf{Run} e \textbf{Air} (locomoção; cada um define a própria animação no inspetor), \textbf{Action} (golpe ou reação a dano em curso), \textbf{Block} (mantém o escudo virado para a câmera e larga a guarda quando um golpe por trás a derruba) e \textbf{Dead} (ignora toda entrada). Cada estado decide a transição seguinte; física, golpes e dano continuam no \texttt{Character} e no \texttt{Attack}, então os quatro aventureiros herdam a máquina sem mudar suas cenas.

\subsection{Rede (co-op)}

Um jogador hospeda (ENet, porta UDP 7777) e os outros entram pelo endereço dele. Cada máquina manda no movimento do próprio aventureiro e envia posição e animação 20 vezes por segundo; o host manda nos mortos-vivos, no elite, no boss e na votação. O dano é aplicado na máquina de quem é dono do alvo. Antes de começar cada fase todos esperam a última máquina carregar, para nenhuma mensagem chegar a uma cena que ainda não existe. Jogar sozinho usa o mesmo código, sem rede. Todo pedido que chega de outra máquina é conferido antes de valer: quem pediu (só o dono do atacante relata o golpe dele), se ele está vivo e perto do que quer mexer (até 35\,m para tiros, o alcance da seta mais o que dá para andar enquanto ela voa) e se o valor é plausível (dano ou cura de no máximo 20, perto do golpe mais forte do jogo). Baú, barril e poção são decididos pelo host, que avisa todas as máquinas juntas; a queda de cada jogador conta uma vez só no placar; depois do início ninguém troca de aventureiro nem de nome; e nomes perdem caracteres de controle e invisíveis.

\subsection{Mapa}

Pedraalta (\texttt{scenes/maps/pedraalta.tscn}) usa os mesmos três GridMaps, com o chão em várias alturas; os terraços, rampas, paredões e a cerca são blocos com colisão, com a textura de pedra (RockTrim) e de tijolo irregular do kit da vila para não sumirem no céu noturno; a decoração grande (entulho, caixotes, chaminés, camas, mesas, portões) tem caixa de colisão do tamanho do modelo, e o navmesh é assado a partir das colisões (\texttt{pedraalta\_navmesh.res}, inclinação máxima de 30 graus; gerado por \texttt{tools/pedraalta/}). Os pontos de grupo levam a zona como metadado, e o sorteio põe no máximo um grupo por zona; mapas sem zona sorteiam como antes. A vila é montada com três GridMaps: \textbf{Chão} (células de 2\,m para ruas e praça), \textbf{Estrutura} (células de 1\,m para pisos, paredes, cantos e escadas das casas) e \textbf{Objetos} (carroças, caixotes, barris e cercas). Telhados, portas, janelas e objetos de cenário ficam fora da grade, como cenas avulsas. As bibliotecas de peças são geradas por um script a partir dos modelos importados. Em volta da vila, quatro paredes invisíveis de 6\,m (\texttt{Limites}) seguram quem sai pelas portas das casas da borda antes do fim do chão. Cada mapa é uma cena própria em \texttt{scenes/maps/} (\texttt{village}, \texttt{grid\_sandbox}, \texttt{public\_square} e \texttt{village\_road}), com os GridMaps \textbf{Chao}, \textbf{Estrutura} e \textbf{Objetos} dentro dela, e é desenhado inteiro pela paleta do GridMap. São três bibliotecas de peças: \textbf{ground} (pisos de pedra e madeira do Medieval Village MegaKit; os pisos do KayKit saíram por serem de cor lisa e destoarem das texturas), \textbf{village\_structure} (paredes, cantos, escadas, telhados e empenas) e \textbf{props} (carroça, caixotes, barris, cercas e tocha de parede, com colisão). Fora da grade ficam só as luzes, os limites invisíveis do mapa, os pontos de sorteio e peças com personagens, como o cadafalso da execução.

\textbf{Floresta de fundo.} Em volta da vila, uma floresta low-poly fecha o horizonte: três árvores modeladas no Blender (pinheiro alto, pinheiro baixo e árvore de copa, de 22 a 32 faces cada) espalhadas com semente fixa entre o quadrado livre de 30\,m de meia largura e o raio de 62\,m, com 2,6\,m entre troncos, desenhadas por \texttt{MultiMeshInstance3D} sem sombra. Embaixo, um chão de terra com textura de ruído; por cima, quatro volumes de neblina rasteira (\texttt{FogVolume}) cobrem o chão entre o mapa e a mata, e as árvores saem da névoa. A floresta é a cena \texttt{scenes/props/forest\_ring.tscn} e serve para qualquer mapa: basta instanciá-la, ajustar no inspetor o quadrado livre, o raio, o espaçamento e a semente e clicar em \emph{Gerar floresta}; o chão e a neblina se ajustam junto. A geração precisa do editor aberto (o modo \texttt{--headless} não guarda as posições das árvores), e a neblina só aparece se o \texttt{Environment} da fase tiver a neblina volumétrica ligada (com densidade global zero, só os volumes contam).

As cutscenes usam o AnimationPlayer do próprio Godot: câmeras, legendas, fades, movimentos e animações dos personagens são trilhas da animação, ajustáveis no painel Animation sem mexer em código. Todas usam o mesmo script (\texttt{AdventurerCutscene}), que só coloca o aventureiro escolhido nos marcadores indicados e, ao terminar, abre a cena definida no inspetor; uma cutscene nova não precisa de código.

\subsection{Testes automatizados}

O projeto usa o gdUnit4 (versão 6.2.1). São NNN testes em SSS suítes, incluindo a mira dos ataques à distância, organizados por área: seleção de aventureiro, assinatura do contrato (nome limpo, salvo e usado na vitória), criação do aventureiro escolhido na partida, cena de execução, um teste específico para cada personagem (Knight, Barbarian, Mage, Rogue e Skeleton Minion), personagem base (vida, morte, buffs), percepção e cérebro dos mortos-vivos (visão, audição, memória, despertar, perseguição, ataque com aviso, busca e navegação contornando paredes), habilidades do botão direito (escudo do Knight, segundo machado do Barbarian), jogo pelo controle (analógicos, gatilhos, filtro por dispositivo, navegação nos menus), formação do elite, movimento do aventureiro (aceleração, andar, correr segurando Espaço, queda, mira pela câmera), comandos de ataque (combo, cancelamento, ataque guardado, golpe interrompido), combate corpo a corpo e à distância, estado da missão, fase da vila (sorteio dos encontros e alcance pela navegação), áudio, baú e poção, Necromante e vitória, configurações, HUD, barra de vida, tela de morte e pausa, máquina de estados do jogador, energia (custos, golpe pesado, dano proporcional, corrida e escudo), votação, histórico de dano e placar, inventário rápido, beber segurando E, troca de habilidade, poção e objetos que quebram, portal e HUD em cruz, funções de rede do lobby, a fase da Pedra Alta (zonas, grupo fácil na entrada e alcance pela navegação), a floresta de fundo, o registro de playtest e o calabouço escondido (peças, escuridão, armadilhas, poção de mana, escadas, barril, espadachim e seus estilos, mímico, contagem de mortos-vivos, saída, travas de classe, selos, alavancas, bilhetes e o alcance de todos os grupos pela navegação). Para rodar: aba gdUnit do editor ou o workflow de testes no GitHub. Além deles, dois testes de ponta a ponta rodam o jogo de verdade a partir do menu, apertando os mesmos botões que o jogador: um contrato solo até a segunda tentativa e uma partida em grupo com dois processos ligados pela rede local (dano entre máquinas, pedido forjado recusado pelo host, reviver pela rede, votação e placar). Rodam com \texttt{tests/e2e/run\_e2e.sh}, também no GitHub.

% ===============================================================
\section{Assets}

O projeto utiliza pacotes gratuitos de modelos 3D e áudio, todos sob licença \href{https://creativecommons.org/publicdomain/zero/1.0/}{CC0} (domínio público, crédito opcional):

\begin{tabularx}{\textwidth}{@{}lX@{}}
\toprule
\textbf{Pacote} & \textbf{Uso no jogo} \\
\midrule
\href{https://kaylousberg.itch.io/kaykit-adventurers}{KayKit Adventurers} (Kay Lousberg) & Aventureiros jogáveis \\
\href{https://kaylousberg.itch.io/kaykit-skeletons}{KayKit Skeletons} (Kay Lousberg) & Mortos-vivos \\
\href{https://kaylousberg.itch.io/kaykit-dungeon-pack}{KayKit Dungeon Pack} (Kay Lousberg) & Props e elementos de cenário (não usado no chão) \\
\href{https://quaternius.itch.io/medieval-village-megakit}{Medieval Village MegaKit} (Quaternius) & Construção da vila \\
\href{https://kenney.nl/assets/rpg-audio}{RPG Audio}, \href{https://kenney.nl/assets/impact-sounds}{Impact Sounds} e \href{https://kenney.nl/assets/interface-sounds}{Interface Sounds} (Kenney) & Efeitos sonoros \\
\href{https://opengameart.org/content/dark-cavern-ambient}{Dark Cavern Ambient} (Paul Wortmann) & Música ambiente da vila \\
\href{https://opengameart.org/content/medieval-fair-loop}{Medieval fair [Loop]} (Woli34) & Música do menu \\
\bottomrule
\end{tabularx}

A fonte dos títulos, \href{https://github.com/NDISCOVER/Cinzel}{Cinzel} (The Cinzel Project Authors), vai embutida no jogo sob a \href{https://openfontlicense.org}{SIL Open Font License 1.1}, que permite uso e redistribuição junto com a licença (\texttt{assets/fonts/OFL.txt}).

\section{Publicação}

O jogo é exportado para Windows e Linux (64 bits) como um executável único, com o conteúdo embutido, e publicado no itch.io. Para jogar em grupo, o host precisa estar na mesma rede local que os outros, ou liberar a porta UDP 7777 no roteador, ou usar uma VPN de jogo (Radmin, Tailscale). A versão para navegador não é usada porque o jogo depende do renderizador Forward+ (sombras, névoa e luzes dinâmicas), que o navegador não suporta. O GitHub Actions gera os executáveis de Windows e Linux em todo pull request para a \texttt{main} (\texttt{build.yml}), publica uma release com os dois a cada merge na \texttt{main} (\texttt{release.yml}) e gera este documento em PDF (\texttt{gdd.yml}). Como o executável não é assinado, o Windows mostra o aviso do SmartScreen na primeira execução (``Mais informações'', ``Executar assim mesmo''); quem hospeda precisa permitir o jogo no firewall, marcando também a rede Pública, que é como redes de faculdade e VPNs costumam aparecer.

O logo é o selo de cera do contrato, partido ao meio, com a lanterna da Companhia da Lanterna guardando um crânio no lugar da chama (\texttt{assets/ui/logo\_seal.svg}). Na interface (ao lado do título no menu e no topo da pausa) ele aparece em uma cor só, dourado sobre o escuro, no tom do título, para não destoar dos botões e da HUD; em cores (cera vermelha e lanterna dourada) ele é o ícone do jogo e do executável; nos tamanhos de 48\,px para baixo o ícone usa o selo inteiro, sem a rachadura, que some nessa escala. O README do projeto traz os controles, como rodar, como testar e os créditos, e os workflows \texttt{tests.yml} e \texttt{e2e.yml} rodam os testes unitários e os de ponta a ponta a cada pull request.

\end{document}
